\IfFileExists{papers/formalized-vs-literature-preamble.tex}{\documentclass[11pt]{article}
\usepackage[utf8]{inputenc}
\usepackage[T1]{fontenc}
\usepackage{amsmath,amssymb,mathtools}
\usepackage{booktabs,tabularx,array}
\usepackage{enumitem}
\usepackage[margin=1in]{geometry}
\usepackage{xcolor}
\usepackage[hidelinks]{hyperref}
\usepackage{microtype}
\setlength{\parindent}{0pt}
\setlength{\parskip}{0.55em}
\newcommand{\Lean}[1]{\texttt{\nolinkurl{#1}}}
\newcommand{\RepoFile}[1]{\texttt{\nolinkurl{#1}}}
\newcommand{\Class}[1]{\textbf{#1}}
\newcommand{\Top}{\mathord{\top}}
\newcommand{\R}{\mathbb{R}}
\newcommand{\ENN}{\mathbb{R}_{\ge 0}^{\infty}}
\newcommand{\KL}{\operatorname{KL}}
\newcommand{\W}{\mathsf{W}}
\newcommand{\statusline}[2]{%
  \begin{center}\fbox{\begin{minipage}{0.92\linewidth}\textbf{Completion class:} #1\\[2pt]
  \textbf{Strongest caveat:} #2\end{minipage}}\end{center}}
\newenvironment{comparison}{\tabularx{\linewidth}{>{\bfseries}p{0.25\linewidth}X}}{\endtabularx}
\newcommand{\snapshot}{\texttt{902f51aa01a737af68f6feb71c83263cc9b1f4d9}}
}{\documentclass[11pt]{article}
\usepackage[utf8]{inputenc}
\usepackage[T1]{fontenc}
\usepackage{amsmath,amssymb,mathtools}
\usepackage{booktabs,tabularx,array}
\usepackage{enumitem}
\usepackage[margin=1in]{geometry}
\usepackage{xcolor}
\usepackage[hidelinks]{hyperref}
\usepackage{microtype}
\setlength{\parindent}{0pt}
\setlength{\parskip}{0.55em}
\newcommand{\Lean}[1]{\texttt{\nolinkurl{#1}}}
\newcommand{\RepoFile}[1]{\texttt{\nolinkurl{#1}}}
\newcommand{\Class}[1]{\textbf{#1}}
\newcommand{\Top}{\mathord{\top}}
\newcommand{\R}{\mathbb{R}}
\newcommand{\ENN}{\mathbb{R}_{\ge 0}^{\infty}}
\newcommand{\KL}{\operatorname{KL}}
\newcommand{\W}{\mathsf{W}}
\newcommand{\statusline}[2]{%
  \begin{center}\fbox{\begin{minipage}{0.92\linewidth}\textbf{Completion class:} #1\\[2pt]
  \textbf{Strongest caveat:} #2\end{minipage}}\end{center}}
\newenvironment{comparison}{\tabularx{\linewidth}{>{\bfseries}p{0.25\linewidth}X}}{\endtabularx}
\newcommand{\snapshot}{\texttt{902f51aa01a737af68f6feb71c83263cc9b1f4d9}}
}
\title{Continuum control-energy/KL identity: formalized result versus the literature}
\author{}
\date{Formalization snapshot \snapshot}
\begin{document}
\maketitle
\statusline{\Class{A/P}: chain-rule decomposition and lower-bound infrastructure are proved; the pathwise Cameron--Martin/Girsanov equality remains an explicit edge}{The central adapted-drift stochastic-exponential theorem is not formalized, so the displayed continuum identity is not yet closed.}
\section{Source targets}
The target identity, appearing in stochastic-control formulations of the Schr\"odinger bridge, is schematically
\[
 \KL(P^{u,\rho_0}\Vert R)
 =\KL(\rho_0\Vert\rho_0^R)
 +\mathbb E_{P^{u,\rho_0}}\int_0^1\frac12\lVert u_t\rVert^2\,dt.
\]
The source chain uses Cameron--Martin for deterministic shifts and Girsanov for adapted drifts.  See
\RepoFile{prose/distilled_literature/Girsanov1960_change_of_measure.tex},
\RepoFile{prose/distilled_literature/CameronMartin1945_wiener_transformations.tex}, and
\RepoFile{prose/distilled_literature/Yeh1978_conditional_wiener_translation.tex}.

\section{Lean declarations}
\Lean{ChenGeorgiouPavon2021.energy_identity} and
\Lean{ChenGeorgiouPavon2021.energy_identity_conditional} are in
\RepoFile{ChenGeorgiouPavon2021/EnergyIdentity.lean}.  Finite and sequence-model instances are in \RepoFile{ChenGeorgiouPavon2021/SequenceGaussian.lean}.

\section{Proved decomposition}
Given a measurable equivalence that separates the initial coordinate and a conditional path, explicit comp-product factorizations of the controlled and reference laws, finite marginal and conditional KL, and the conditional Cameron--Martin identity, Lean proves the full KL decomposition by:
\begin{enumerate}[leftmargin=2em]
\item invariance of KL under measurable equivalence;
\item the kernel/product KL chain rule;
\item substitution of the conditional energy identity.
\end{enumerate}
A conditional version rewrites the conditional KL as the nominal average of pointwise kernel KL.  Projection/data-processing lemmas also prove the It\^o-free lower-bound direction once projected KL values converge to the energy.

\section{Comparison}
\begin{comparison}
KL chain rule & Proved.\\
Conditional KL integral & Proved under countable-generation assumptions.\\
Finite Gaussian energy atom & Proved.\\
Iid sequence deterministic shift & Proved completely.\\
Continuous deterministic path shift & Not yet transported to a canonical continuous-path law.\\
Adapted drift/Girsanov & Open.\\
It\^o integral/SDE path laws & Open in the current dependency stack.\\
\end{comparison}

\section{Nonclaims}
The theorem named \Lean{energy_identity} is an assembly theorem over a genuine nonvacuous interface.  It must not be described as a completed proof of Girsanov or of the continuum DRSB energy identity.
\end{document}
